% ============================================================
%  ROADMAP PERSONAL
%  Rosh Guadiana — Uso exclusivamente privado
%  Compilar con: pdflatex (dos pasadas)
%  Versi\'on: 8 Sesiones (Destilada con Contexto Hist\'orico y Bibliograf\'ia)
% ============================================================

\documentclass[11pt, letterpaper]{article}

\usepackage[utf8]{inputenc}
\usepackage{geometry}
\geometry{top=2.2cm, bottom=2.5cm, left=2.8cm, right=2.8cm, headheight=22pt}
\usepackage{amsmath, amssymb}
\usepackage{booktabs, array}
\usepackage[dvipsnames]{xcolor}
\usepackage{fancyhdr}
\usepackage{titlesec}
\usepackage[colorlinks=true, linkcolor=roshblue, urlcolor=roshblue]{hyperref}
\usepackage{tcolorbox}
\tcbuselibrary{skins, breakable}
\usepackage{enumitem}
\usepackage{parskip}
\usepackage{graphicx}
\usepackage{mdframed}
\usepackage{calc}

% --- Colors ---
\definecolor{roshblue}{RGB}{10, 60, 120}
\definecolor{roshgreen}{RGB}{20, 100, 50}
\definecolor{roshorange}{RGB}{180, 80, 0}
\definecolor{roshred}{RGB}{150, 20, 20}
\definecolor{lightgray}{RGB}{245, 245, 245}
\definecolor{warningyellow}{RGB}{255, 248, 220}
\definecolor{warningborder}{RGB}{200, 150, 0}
\definecolor{okgreen}{RGB}{230, 245, 230}
\definecolor{okborder}{RGB}{20, 120, 50}
\definecolor{historypurple}{RGB}{90, 30, 110}
\definecolor{refgray}{RGB}{240, 240, 240}

% --- Checkbox command ---
\newcommand{\cbox}{\makebox[1.2em]{$\square$}}
\newcommand{\cboxdone}{\makebox[1.2em]{\textcolor{roshgreen}{$\checkmark$}}}

% --- Header ---
\pagestyle{fancy}
\fancyhf{}
\fancyhead[L]{\small\textcolor{roshblue}{\textsc{Roadmap Personal}}}
\fancyhead[R]{\small\textcolor{roshblue}{Rosh Guadiana --- Privado}}
\fancyfoot[C]{\small\textcolor{gray}{\thepage}}
\renewcommand{\headrulewidth}{0.4pt}
\renewcommand{\headrule}{\hbox to\headwidth{%
  \color{roshblue}\leaders\hrule height \headrulewidth\hfill}}

% --- Sections ---
\titleformat{\section}
  {\large\bfseries\color{roshblue}}
  {\thesection.}{0.6em}{}
  [\vspace{-0.3em}\textcolor{roshblue}{\rule{\linewidth}{0.5pt}}]
\titleformat{\subsection}
  {\normalsize\bfseries\color{roshblue!80}}{\thesubsection}{0.5em}{}

% --- Boxes ---
\newtcolorbox{verificacion}[1]{
  colback=lightgray, colframe=roshblue,
  fonttitle=\bfseries\small,
  title=Rigor Matem\'atico: #1,
  left=5pt, right=5pt, top=4pt, bottom=4pt,
  boxrule=0.6pt, breakable}
  
\newtcolorbox{historia}[1]{
  colback=white, colframe=historypurple,
  fonttitle=\bfseries\small,
  title=Rigor Hist\'orico y Biogr\'afico: #1,
  left=5pt, right=5pt, top=4pt, bottom=4pt,
  boxrule=0.6pt, breakable}

\newtcolorbox{referencias}{
  colback=refgray, colframe=gray,
  fonttitle=\bfseries\small\color{darkgray},
  title=Fuentes y Referencias,
  left=5pt, right=5pt, top=4pt, bottom=4pt,
  boxrule=0.4pt, breakable}

\newtcolorbox{advertencia}{
  colback=warningyellow, colframe=warningborder,
  fonttitle=\bfseries\small\color{warningborder},
  title=Atenci\'on,
  left=5pt, right=5pt, top=4pt, bottom=4pt,
  boxrule=0.6pt, breakable}

\newtcolorbox{notapersonal}{
  colback=okgreen, colframe=okborder,
  fonttitle=\bfseries\small,
  title=Nota Personal,
  left=5pt, right=5pt, top=4pt, bottom=4pt,
  boxrule=0.5pt, breakable}

% --- List styles ---
\newlist{checklist}{itemize}{2}
\setlist[checklist]{
  label=\cbox,
  leftmargin=2em,
  itemsep=4pt,
  topsep=4pt}

\setlist[itemize]{leftmargin=1.4em, itemsep=2pt, topsep=3pt}
\setlength{\parskip}{5pt}

% ============================================================
\begin{document}

\begin{titlepage}
\begin{center}
\vspace*{1cm}
{\footnotesize\textsc{\textcolor{gray}{Documento de uso exclusivamente personal}}}
\vspace{0.5cm}
\textcolor{roshblue}{\rule{\linewidth}{1pt}}
\vspace{0.4cm}
{\LARGE\bfseries\color{roshblue} Roadmap Personal}\\[0.4em]
{\large\itshape\color{gray} Dominio Matem\'atico, Contexto Hist\'orico y Plan de Estudio}
\vspace{0.3cm}
\textcolor{roshblue}{\rule{\linewidth}{1pt}}
\vspace{1.5cm}

\begin{tcolorbox}[colback=warningyellow, colframe=warningborder, boxrule=0.6pt]
Este documento contiene lo que debes saber \textbf{sin mirar ning\'un libro}. Si no puedes derivar la f\'ormula o contar la an\'ecdota hist\'orica de memoria, a\'un no est\'as listo para esa sesi\'on.
\end{tcolorbox}
\end{center}
\end{titlepage}

\newpage

% ============================================================
%  PLAN DE ESTUDIO INTENSIVO (10 DÍAS)
% ============================================================
\section{Plan de Estudio Intensivo (10 D\'ias)}

\begin{center}
\renewcommand{\arraystretch}{1.5}
\begin{tabular}{clp{9.5cm}}
\toprule
\textbf{D\'ia} & \textbf{Foco} & \textbf{Objetivo Diario} \\
\midrule
1 & Conjuntos y Cantor & Halmos Cap. 1-8. Memorizar estructura de la prueba diagonal. Investigar conflicto Cantor-Kronecker. \\
2 & L\'ogica Booleana & Derivaciones de completitud funcional (NAND). Leer historia de Boole y su intenci\'on cognitiva. \\
3 & L\'ogica 1er Orden & Sintaxis, sem\'antica y modelos (Enderton/Deaño). Entender K\"onigsberg 1930 y el choque Hilbert-G\"odel. \\
4 & G\"odel Formal & Nagel \& Newman Cap. 1-7. Trazar en papel la aritmetizaci\'on de la sintaxis y la f\'ormula G. \\
5 & Turing y Conciencia & Sipser Cap. 4 y 6. Demostrar el Problema de la Parada. Leer biograf\'ia de Turing y \textit{The Emperor's New Mind} Cap. 10. \\
6 & El Isomorfismo & Shannon (1938). Entender f\'isicamente por qu\'e la linealidad destruye la se\~nal (ruido y atenuaci\'on). \\
7 & Fermi-Dirac & Streetman Cap. 3. Derivar la funci\'on de distribuci\'on desde combinatoria/exclusi\'on de Pauli. \\
8 & La Uni\'on PN & Streetman Cap. 5. Derivaci\'on del potencial built-in y la ecuaci\'on de Shockley. Historia de Bell Labs (1947). \\
9 & F\'isica del MOSFET & Sedra Cap. 5. Derivaci\'on de corriente de canal largo ($I_D$) y pinzamiento (pinch-off). Historia de Kahng/Atalla (1959). \\
10 & CMOS L\'ogico & Sedra Cap. 14. Trazar redes Pull-Up/Pull-Down en papel. S\'intesis final del curso. Ensayo general mental. \\
\bottomrule
\end{tabular}
\end{center}

\begin{notapersonal}
No intentes leer los libros de texto completos de portada a contraportada. Ve directamente a los cap\'itulos y secciones exactas marcadas en las referencias de cada sesi\'on. Estudia con pluma en mano: si no lo puedes escribir en una hoja en blanco, no lo has memorizado.
\end{notapersonal}

\newpage

% ============================================================
%  SESIÓN 1
% ============================================================
\section{Sesi\'on 1 --- Fundamentos y el Infinito}

\begin{verificacion}{Sesi\'on 1}
\begin{checklist}
  \item Define formalmente $A\cup B$, $A\cap B$, $A\setminus B$ y el conjunto potencia $\mathcal{P}(A)$.
  \item Demuestra en papel que $|\mathcal{P}(A)| = 2^{|A|}$ para $A=\{a,b,c\}$.
  \item Escribe la paradoja de Russell formalmente: $R=\{x \mid x\notin x\}$. Demuestra el colapso.
  \item Construye la biyecci\'on entre $\mathbb{N}$ y $\mathbb{Z}$ expl\'icitamente.
  \item Dibuja la matriz diagonal de Cantor y construye el n\'umero $d$ que escapa de la lista.
\end{checklist}
\end{verificacion}

\begin{historia}{Sesi\'on 1}
\begin{checklist}
  \item \textbf{La carta de Russell a Frege (1902):} Investiga c\'omo Gottlob Frege estaba a punto de publicar el segundo volumen de su obra magna fundando las matem\'aticas enteras, cuando recibi\'o una carta de Bertrand Russell con la paradoja, destruyendo el trabajo de su vida.
  \item \textbf{Cantor y Kronecker:} Investiga la relaci\'on t\'oxica entre Leopold Kronecker (el finitista que dec\'ia \textit{"Dios hizo los n\'umeros enteros, el resto es obra del hombre"}) y Cantor. C\'omo Kronecker bloque\'o la carrera de Cantor y lo empuj\'o a la depresi\'on.
\end{checklist}
\end{historia}

\begin{referencias}
\textbf{Para Matem\'atica:}
\begin{itemize}
    \item Halmos, P. R. --- \textit{Naive Set Theory} (Caps. 1--8).
    \item Enderton, H. B. --- \textit{Elements of Set Theory} (Cap. 1).
\end{itemize}
\textbf{Para Historia:}
\begin{itemize}
    \item Aczel, A. D. --- \textit{The Mystery of the Aleph} (Historia de Cantor, su depresi\'on y Kronecker).
    \item Doxiadis, A. \& Papadimitriou, C. H. --- \textit{Logicomix} (Excelente para el drama humano de Russell, Frege y la paradoja).
\end{itemize}
\end{referencias}

\newpage

% ============================================================
%  SESIÓN 2
% ============================================================
\section{Sesi\'on 2 --- L\'ogica Proposicional}

\begin{verificacion}{Sesi\'on 2}
\begin{checklist}
  \item Define FBF de manera puramente inductiva.
  \item Demuestra que Modus Ponens $\{p\to q, p\} \models q$ es una tautolog\'ia usando tabla de verdad.
  \item \textbf{Crucial:} Demuestra algebraicamente que $\{\text{NAND}\}$ es funcionalmente completo, construyendo $\neg p$ y $p\land q$ solo con NANDs.
  \item Escribe una prueba de derivaci\'on formal simple para separar $\vdash$ (sintaxis) de $\models$ (sem\'antica).
\end{checklist}
\end{verificacion}

\begin{historia}{Sesi\'on 2}
\begin{checklist}
  \item \textbf{George Boole y \textit{The Laws of Thought} (1854):} Investiga cu\'al era la verdadera motivaci\'on de Boole. No quer\'ia hacer electr\'onica (no exist\'ia); quer\'ia modelar matem\'aticamente c\'omo funciona el pensamiento de la mente humana.
  \item \textbf{La tragedia de Boole:} C\'omo muri\'o de neumon\'ia despu\'es de caminar bajo la lluvia y que su esposa (creyendo que el fr\'io curaba los efectos del fr\'io) lo tratara con cubetas de agua helada.
\end{checklist}
\end{historia}

\begin{referencias}
\textbf{Para Matem\'atica:}
\begin{itemize}
    \item Deaño, A. --- \textit{Introducci\'on a la l\'ogica formal} (Partes I y II).
    \item Enderton, H. B. --- \textit{A Mathematical Introduction to Logic} (Cap. 1).
\end{itemize}
\textbf{Para Historia:}
\begin{itemize}
    \item MacTutor History of Mathematics Archive (St Andrews University) --- \textit{Biograf\'ia de George Boole}.
    \item Gleick, J. --- \textit{The Information} (Cap\'itulo sobre Boole y su intenci\'on de mecanizar el pensamiento).
\end{itemize}
\end{referencias}

\newpage

% ============================================================
%  SESIÓN 3
% ============================================================
\section{Sesi\'on 3 --- L\'ogica de Primer Orden (FOL)}

\begin{verificacion}{Sesi\'on 3}
\begin{checklist}
  \item Define firma l\'ogica $\mathcal{L}=(\mathcal{F},\mathcal{P},\text{ar})$ con ejemplos.
  \item Distingue matem\'aticamente entre variables libres y ligadas en una FBF con cuantificadores.
  \item Define el concepto de Estructura $\mathfrak{A}$ (dominio e interpretaci\'on).
  \item Enuncia rigurosamente el Teorema de Completitud de G\"odel (1929) ($\Gamma\models\varphi \Leftrightarrow \Gamma\vdash\varphi$).
\end{checklist}
\end{verificacion}

\begin{historia}{Sesi\'on 3}
\begin{checklist}
  \item \textbf{El Programa de Hilbert (d\'ecada de 1920):} Investiga el sue\~no de David Hilbert de mecanizar todas las matem\'aticas para probar que estaban libres de contradicci\'on. (Aprende su lema: \textit{"Wir m\"ussen wissen. Wir werden wissen"}).
  \item \textbf{El Congreso de K\"onigsberg (1930):} El momento en que Hilbert daba su famoso discurso optimista, mientras en una mesa redonda sat\'elite en la misma ciudad, G\"odel (de 24 a\~nos) anunciaba que el sue\~no de Hilbert estaba matem\'aticamente muerto.
\end{checklist}
\end{historia}

\begin{referencias}
\textbf{Para Matem\'atica:}
\begin{itemize}
    \item Enderton, H. B. --- \textit{A Mathematical Introduction to Logic} (Cap. 2).
    \item Marker, D. --- \textit{Model Theory: An Introduction} (Cap. 1, p\'aginas iniciales para rigor de estructuras).
\end{itemize}
\textbf{Para Historia:}
\begin{itemize}
    \item Dawson, J. W. --- \textit{Logical Dilemmas: The Life and Work of Kurt G\"odel} (Contexto del congreso de K\"onigsberg).
    \item Reid, C. --- \textit{Hilbert} (Para entender la ambici\'on del Programa de Hilbert).
\end{itemize}
\end{referencias}

\newpage

% ============================================================
%  SESIÓN 4
% ============================================================
\section{Sesi\'on 4 --- Metamatem\'atica y Computabilidad}

\begin{verificacion}{Sesi\'on 4}
\begin{checklist}
  \item Explica la aritmetizaci\'on de la sintaxis (N\'umeros de G\"odel).
  \item Escribe la estructura l\'ogica de la oraci\'on $G$ ("Yo no soy demostrable en este sistema").
  \item Dibuja la misma matriz diagonal para el Problema de la Parada (M\'aquina $H$ vs M\'aquina $D$).
  \item Enuncia con precisi\'on el 1er y 2do Teorema de Incompletitud.
\end{checklist}
\end{verificacion}

\begin{historia}{Sesi\'on 4}
\begin{checklist}
  \item \textbf{John von Neumann y G\"odel:} Investiga c\'omo, durante el anuncio de G\"odel en 1930, von Neumann fue el \'unico en la audiencia que entendi\'o inmediatamente las consecuencias destructivas del teorema, arrinconando a G\"odel despu\'es de la charla.
  \item \textbf{Alan Turing (1936):} Investiga su \textit{paper} "On Computable Numbers". C\'omo Turing invent\'o la computadora conceptualmente (la m\'aquina autom\'atica) s\'olo para probar que el problema de la decisi\'on de Hilbert (\textit{Entscheidungsproblem}) no ten\'ia soluci\'on.
\end{checklist}
\end{historia}

\begin{referencias}
\textbf{Para Matem\'atica:}
\begin{itemize}
    \item Sipser, M. --- \textit{Introduction to the Theory of Computation} (Cap. 4 y 6, problema de la parada).
    \item Nagel, E. \& Newman, J. R. --- \textit{G\"odel's Proof} (Todo el libro es oro para la prueba paso a paso).
\end{itemize}
\textbf{Para Historia:}
\begin{itemize}
    \item Turing, A. M. (1936) --- \textit{On Computable Numbers, with an Application to the Entscheidungsproblem} (Fuente primaria).
    \item Hodges, A. --- \textit{Alan Turing: The Enigma} (Mejor biograf\'ia de Turing).
\end{itemize}
\end{referencias}

\newpage

% ============================================================
%  SESIÓN 5
% ============================================================
\section{Sesi\'on 5 --- Di\'alogo Perip\'at\'etico}

\begin{historia}{Sesi\'on 5 (Aire Libre)}
\begin{checklist}
  \item \textbf{Narrativa entrelazada:} Debes poder contar, sin leer y con ritmo pausado, c\'omo las tres mentes que definieron los l\'imites de la raz\'on (Cantor, G\"odel, Turing) tuvieron finales tr\'agicos.
  \item \textbf{G\"odel:} Su paranoia, su amistad con Einstein en Princeton (Einstein dec\'ia que iba a la oficina \textit{"solo para tener el privilegio de caminar a casa junto a G\"odel"}), y su muerte pesando 29 kg por inanici\'on por miedo a ser envenenado tras la hospitalizaci\'on de su esposa.
  \item \textbf{Turing:} Su trabajo fundamental en Bletchley Park rompiendo la m\'aquina Enigma nazi, su juicio por indecencia (homosexualidad) por el mismo gobierno al que salv\'o, la castraci\'on qu\'imica y la manzana con cianuro.
\end{checklist}
\end{historia}

\begin{referencias}
\textbf{Para el Di\'alogo Filos\'ofico:}
\begin{itemize}
    \item Hofstadter, D. R. --- \textit{G\"odel, Escher, Bach} (Especialmente Caps. IX, XV, y XX sobre inteligencia artificial y bucles extra\~nos).
    \item Penrose, R. --- \textit{The Emperor's New Mind} (Cap. 10: Argumento de que la conciencia es no-computable).
\end{itemize}
\textbf{Para Biograf\'ias Tr\'agicas:}
\begin{itemize}
    \item Goldstein, R. --- \textit{Incompleteness: The Proof and Paradox of Kurt G\"odel} (Los \'ultimos a\~nos de G\"odel en Princeton).
    \item Hodges, A. --- \textit{Alan Turing: The Enigma} (El juicio y final de Turing).
\end{itemize}
\end{referencias}

\newpage

% ============================================================
%  SESIÓN 6
% ============================================================
\section{Sesi\'on 6 --- El Puente (Boole y Shannon)}

\begin{verificacion}{Sesi\'on 6}
\begin{checklist}
  \item Explica la equivalencia uno a uno entre funciones l\'ogicas abstractas y redes de interruptores (abierto/cerrado).
  \item Justifica termodin\'amicamente por qu\'e la linealidad ($V_\text{out}=\alpha V_\text{in}$) destruye la precisi\'on de la se\~nal (amplificaci\'on de ruido sin regeneraci\'on).
  \item Define asimetr\'ia/no-linealidad como requisito f\'isico absoluto para la retenci\'on de un bit.
\end{checklist}
\end{verificacion}

\begin{historia}{Sesi\'on 6}
\begin{checklist}
  \item \textbf{Claude Shannon (1938):} Investiga su tesis de maestr\'ia en MIT. C\'omo tom\'o una clase electiva de filosof\'ia sobre \'algebra de Boole (considerada in\'util en ese momento) y not\'o que la l\'ogica simb\'olica encajaba exactamente con el enrutamiento complejo de los rel\'es telef\'onicos mec\'anicos que operaba. Es considerada la tesis de maestr\'ia m\'as influyente de la historia.
\end{checklist}
\end{historia}

\begin{referencias}
\textbf{Para Matem\'atica/F\'isica:}
\begin{itemize}
    \item Shannon, C. E. (1938) --- \textit{A Symbolic Analysis of Relay and Switching Circuits} (Trans. AIEE).
    \item Tanenbaum, A. S. --- \textit{Structured Computer Organization} (Cap. 1, de la l\'ogica al circuito).
\end{itemize}
\textbf{Para Historia:}
\begin{itemize}
    \item Soni, J. \& Goodman, R. --- \textit{A Mind at Play: How Claude Shannon Invented the Information Age} (La historia de su tesis en MIT).
\end{itemize}
\end{referencias}

\newpage

% ============================================================
%  SESIÓN 7
% ============================================================
\section{Sesi\'on 7 --- Termodin\'amica L\'ogica ($E_g$ y Uni\'on PN)}

\begin{verificacion}{Sesi\'on 7}
\begin{checklist}
  \item Deriva la funci\'on de Fermi-Dirac asumiendo el principio de exclusi\'on de Pauli (f\'isica estad\'istica).
  \item Explica $E_g$ como barrera l\'ogica. ¿Qu\'e pasa con $n_i(T)$ si $k_BT \approx E_g$?
  \item Deriva por qu\'e la difusi\'on de portadores genera el potencial built-in ($V_0$) en la uni\'on PN.
  \item Demuestra matem\'aticamente c\'omo la polarizaci\'on de la uni\'on arroja la ecuaci\'on asim\'etrica de Shockley ($I_D \propto e^{V/V_T}$).
\end{checklist}
\end{verificacion}

\begin{historia}{Sesi\'on 7}
\begin{checklist}
  \item \textbf{Bell Labs (1947):} La invenci\'on del primer transistor (de punto de contacto). Investiga la din\'amica t\'oxica entre William Shockley y su equipo (Bardeen y Brattain).
  \item \textbf{El mes m\'agico:} C\'omo Bardeen y Brattain lograron el efecto el\'ectrico sin Shockley. C\'omo Shockley, furioso por haber sido excluido de la patente, se encerr\'o en una habitaci\'on de hotel en Chicago por semanas y desarroll\'o la teor\'ia completa del transistor de uni\'on bipolar, s\'olo para probar que era superior a ellos.
\end{checklist}
\end{historia}

\begin{referencias}
\textbf{Para Matem\'atica/F\'isica:}
\begin{itemize}
    \item Streetman, B. G. \& Banerjee, S. K. --- \textit{Solid State Electronic Devices} (Caps. 3 y 5. La demostraci\'on de Fermi-Dirac est\'a en los ap\'endices termodin\'amicos del Cap. 3).
    \item Sedra, A. S. \& Smith, K. C. --- \textit{Microelectronic Circuits} (Cap. 4, l\'ogica del diodo).
\end{itemize}
\textbf{Para Historia:}
\begin{itemize}
    \item Isaacson, W. --- \textit{The Innovators} (Cap. 4: El transistor. Excelente detalle de la pelea de egos en Bell Labs).
    \item Riordan, M. \& Hoddeson, L. --- \textit{Crystal Fire: The Invention of the Transistor and the Birth of the Information Age}.
\end{itemize}
\end{referencias}

\newpage

% ============================================================
%  SESIÓN 8
% ============================================================
\section{Sesi\'on 8 --- El Modus Ponens F\'isico (MOSFET y CMOS)}

\begin{verificacion}{Sesi\'on 8}
\begin{checklist}
  \item Dibuja la secci\'on transversal de un MOSFET. Explica c\'omo el $V_{GS}$ genera inversi\'on por efecto de campo.
  \item Deriva la ecuaci\'on de la corriente de saturaci\'on $I_D$ (pinch-off).
  \item Dise\~na la red CMOS (Pull-Up y Pull-Down) para un inversor.
  \item Dise\~na la topolog\'ia CMOS transistorial para la compuerta NAND universal.
  \item \textbf{El arco final:} Conecta axiom\'aticamente en una sola narraci\'on matem\'atica el s\'imbolo l\'ogico $\neg(A \land B)$ con la carga cu\'antica de la red CMOS.
\end{checklist}
\end{verificacion}

\begin{historia}{Sesi\'on 8}
\begin{checklist}
  \item \textbf{Dawon Kahng y Martin Atalla (1959):} Investiga c\'omo inventaron el MOSFET (el transistor de efecto de campo, que hoy forma el 99.9\% de la econom\'ia digital) en Bell Labs.
  \item \textbf{La Ceguera Corporativa:} C\'omo Bell Labs engavet\'o el MOSFET por a\~nos porque la gerencia estaba obsesionada con perfeccionar el transistor bipolar de Shockley. Tuvieron que ser compa\~n\'ias externas (Fairchild, Intel) quienes apostaran por la l\'ogica MOS.
\end{checklist}
\end{historia}

\begin{referencias}
\textbf{Para Matem\'atica/F\'isica:}
\begin{itemize}
    \item Sedra, A. S. \& Smith, K. C. --- \textit{Microelectronic Circuits} (Caps. 5 y 14. Aqu\'i est\'a el dise\~no exacto del inversor y NAND en CMOS).
    \item Weste, N. H. E. \& Harris, D. --- \textit{CMOS VLSI Design} (Cap. 1, topolog\'ias l\'ogicas y consumo nulo de energ\'ia est\'atica).
\end{itemize}
\textbf{Para Historia:}
\begin{itemize}
    \item Bassett, R. K. --- \textit{To the Digital Age: Research Labs, Start-up Companies, and the Rise of MOS Technology} (La ignorancia de Bell Labs hacia el MOSFET).
    \item Lojek, B. --- \textit{History of Semiconductor Engineering} (Detalles del trabajo de Atalla y Kahng).
\end{itemize}
\end{referencias}

\end{document}
